% NS-52. Exact finite identities and a conditional obstruction; no RH assumption.
% Standalone fragment: requires amsmath, amssymb, amsthm, hyperref and
% proposition/proof environments. No manuscript version is assigned here.
\subsection*{NS-52: paired heat-zero transport and its first obstruction}

\paragraph{Normalization and claim scope.}
Use the conventional de Bruijn--Newman family
\[
 H_t(z)=\int_0^\infty e^{tu^2}\Phi(u)\cos(zu)\,du,
 \qquad H_0(z)=\tfrac18\xi(\tfrac12+iz/2),
\]
as in Polymath15, equations (1)--(4) and Proposition 3.1
(\url{https://arxiv.org/html/1904.12438#S3}).
Put $Z_t(w)=8H_t(2w)$ and $\kappa=1/4$. Then
\begin{equation}\label{eq:ns52-normalization}
 Z_0(w)=\Xi(w):=\xi(\tfrac12+iw),\qquad
 \partial_t Z_t=-\kappa\partial_w^2 Z_t.
\end{equation}
Thus the heat coefficient in the ordinary zeta-ordinate coordinate is
$1/4$, not $1$. All zero sums below count multiplicity. No assertion
about the locations of the zeros of $Z_0$ is assumed.

For $f,g\in C_c^\infty(-a,a)$ define the nonunitary transform and
its analytic sesquilinear product by
\[
 F_f(w)=\int_{-a}^a f(x)e^{-iwx}\,dx,\qquad
 F_g^*(w)=\overline{F_g(\overline w)},\qquad
 \Psi_{f,g}(w)=F_f(w)F_g^*(w).
\]
The second factor is an entire function of $w$. Replacing it by
$\overline{F_g(w)}$ off the real axis would give the wrong form.
With $\mathcal Z_t$ the zero multiset of $Z_t$, put
\begin{equation}\label{eq:ns52-zero-form}
 q_{a,t}(f,g)=\sum_{w\in\mathcal Z_t}\Psi_{f,g}(w).
\end{equation}
At $t=0$, $\rho=1/2+iw$ and
$\rho^\sharp=1-\overline\rho=1/2+i\overline w$ identify this
exactly with the manuscript's first-slot-linear Weil zero formula:
$q_{a,0}(f,g)=q_a(f,g)$. There is no $2\pi$ factor with the
displayed nonunitary transform. Using the unitary transform instead
would put $2\pi$ in front of the sum. The factor $8$ in $Z_t$ does
not change any zero or its multiplicity. The completed function has
no poles; no additional pole term is to be added to this zero-side
definition. This is not an assertion that the undeformed geometric
prime and pole terms can be heat-evolved separately.

For completeness, the sum in \eqref{eq:ns52-zero-form} is absolutely
convergent on this smooth core for every $t\ge0$, locally uniformly
in $t$. The de Bruijn strip theorem (Polymath15, Theorem 3.2) and
the unconditional initial critical strip give
$|\operatorname{Im} w|\le1/2$. For $0\le t\le T$, Gaussian
majorization of the super-exponentially decreasing $\Phi$ gives
$|Z_t(w)|\le C_T e^{C_T|w|^2}$ and $Z_t(0)\ge Z_0(0)>0$.
Jensen's formula therefore bounds the zero count in $|w|\le R$
by $C_T(1+R)^2$, uniformly in $t$. Repeated integration by parts
gives arbitrary inverse powers of $|\operatorname{Re}w|$ for both
test transforms on the fixed strip. This proves convergence and
uniformly small tails. Local contour integrals count any finite
cluster continuously through collisions, so the complete form is
continuous in $t$ on the smooth core. This argument does not justify
termwise differentiation of the infinite sum or claim a closed
deformed form on the original complete form domain.
If all zeros at a fixed $t_0$ are real, then
\[
 q_{a,t_0}[f]=\sum_{w\in\mathcal Z_{t_0}}|F_f(w)|^2\ge0.
\]
For example the classical de Bruijn theorem supplies such $t_0\ge1/2$.

\begin{proposition}[Finite-cluster transport, including a collision]
\label{prop:ns52-cluster}
Let $Z(t,w)$ be holomorphic near $J\times\overline\Omega$, where
$J$ is an interval of real times and $\Omega$ is a bounded domain
with piecewise smooth positively oriented boundary $\Gamma$.
Assume $Z$ has no zero on $\Gamma$ for $t\in J$ and satisfies
$\partial_t Z=-\kappa Z_{ww}$ with constant $\kappa>0$.
For an entire $\Psi$, define the finite sum $S_\Gamma(t)$ of
$\Psi$ over the zeros in $\Omega$, counted with multiplicity.
Then $S_\Gamma$ is differentiable through all internal collisions and
\begin{align}
 S_\Gamma(t)&=\frac1{2\pi i}\int_\Gamma
                   \Psi(w)\frac{Z_w}{Z}(t,w)\,dw,
                   \label{eq:ns52-contour}\\
 S_\Gamma'(t)&=\frac{\kappa}{2\pi i}\int_\Gamma
                   \Psi'(w)\frac{Z_{ww}}{Z}(t,w)\,dw.
                   \label{eq:ns52-contour-derivative}
\end{align}
At a time with $m$ simple cluster zeros $w_1,\ldots,w_m$, set
$A(w)=Z(t,w)/\prod_{j=1}^m(w-w_j)$ and $b=A'/A$ locally in
the cluster. Then
\begin{equation}\label{eq:ns52-divided-difference}
 S_\Gamma'
 =2\kappa\sum_{1\le j<k\le m}
       \frac{\Psi'(w_j)-\Psi'(w_k)}{w_j-w_k}
   +2\kappa\sum_{j=1}^m b(w_j)\Psi'(w_j).
\end{equation}
If the cluster consists at a particular time of one zero $c$ of
order $m$, write $Z(t,w)=(w-c)^m A(w)$ with $A(c)\ne0$. Its
derivative at that time is exactly
\begin{equation}\label{eq:ns52-multiple}
 S_\Gamma'=
 \kappa m(m-1)\Psi''(c)
       +2\kappa m\frac{A'(c)}{A(c)}\Psi'(c).
\end{equation}
In particular, a double zero contributes
$2\kappa\Psi''(c)+4\kappa(A'/A)(c)\Psi'(c)$.
\end{proposition}
\begin{proof}
The argument principle proves \eqref{eq:ns52-contour}. The nonvanishing
boundary allows differentiation under its compact integral. Since
\[
 \partial_t(Z_w/Z)=-\kappa\partial_w(Z_{ww}/Z),
\]
integration by parts around $\Gamma$ gives
\eqref{eq:ns52-contour-derivative}. At a simple zero,
$w_j'=\kappa Z_{ww}(w_j)/Z_w(w_j)$, and factorization gives
\[
 w_j'=2\kappa\sum_{k\ne j}\frac1{w_j-w_k}
                         +2\kappa b(w_j).
\]
Pair the $j,k$ and $k,j$ contributions to $\sum_jw_j'\Psi'(w_j)$.
This proves \eqref{eq:ns52-divided-difference}. At an order-$m$ zero,
\[
 \frac{Z_{ww}}Z=\frac{m(m-1)}{(w-c)^2}
       +\frac{2m}{w-c}\frac{A'}A+\frac{A''}A.
\]
The residue after multiplying by $\Psi'$ is precisely
\eqref{eq:ns52-multiple}. No labeling or derivative of a colliding
individual root is needed.
\end{proof}

\paragraph{Complete symmetries and the term that must not be omitted.}
For real $t$, $Z_t$ is real entire and even. Choose the union of
cluster contours invariant under conjugation and $w\mapsto-w$,
counting each distinct zero once with its actual multiplicity.
Then the sum defines a Hermitian sesquilinear contribution and the
same identities apply. A double collision at a nonzero real $c$
has a reflected double collision at $-c$; both must be included
in the symmetry-complete contribution. At $c=0$ one would count
one cluster, not duplicate it (the actual $H_t$ has no imaginary-axis
zero). There is no RH assumption here.

For an isolated pair $w_\pm=c\pm d$, the internal term is
$2\kappa[\Psi'(w_+)-\Psi'(w_-)]/(w_+-w_-)$.
It converges to $2\kappa\Psi''(c)$ as $d\to0$.
The two divergent individual velocities therefore cancel exactly.
The remaining $b=A'/A$ term retains the interaction with the rest
of the entire function; it does not vanish by symmetry. The local
analytic factor suffices for the proof, so no additional rearrangement
of an infinite velocity sum is used here. For a full isolated double
cluster the limiting value is the double-zero formula above, not
merely $2\kappa\Psi''(c)$.

\begin{proposition}[The elementary estimate still grows with support]
\label{prop:ns52-support}
For $f,g\in L^2(-a,a)$ and $n\ge0$,
\begin{equation}\label{eq:ns52-pw}
 |\Psi_{f,g}^{(n)}(w)|
 \le 2a(2a)^n e^{2a|\operatorname{Im}w|}\|f\|_2\|g\|_2.
\end{equation}
Consequently, at a double collision with
$|\operatorname{Im}c|\le Y$ and $|(A'/A)(c)|\le B$, the bound
\begin{equation}\label{eq:ns52-support-cost}
 |S_\Gamma'|
 \le16\kappa(a^3+Ba^2)e^{2aY}\|f\|_2\|g\|_2
\end{equation}
holds. The $a^3$ dependence of a single real collision contribution
cannot be replaced by a constant for all compact smooth tests.
Specifically, suppose a double collision occurs at real $c$, and
choose a real odd $h\in C_c^\infty(-1,1)$ with $\|h\|_2=1$ and
$M_1=\int u h(u)\,du\ne0$. For
$f_a(x)=a^{-1/2}h(x/a)e^{icx}$ one has exactly
\begin{equation}\label{eq:ns52-collision-witness}
 \|f_a\|_2=1,\qquad
 S_\Gamma'[f_a]=4\kappa a^3 M_1^2.
\end{equation}
The occurrence of such a collision in the nonnegative-time arithmetic
family is not asserted.
\end{proposition}
\begin{proof}
Cauchy--Schwarz gives
$|F_f^{(j)}(w)|\le a^j\sqrt{2a}e^{a|\operatorname{Im}w|}\|f\|_2$.
Leibniz's formula proves \eqref{eq:ns52-pw}; use $n=1,2$ in
\eqref{eq:ns52-multiple} to obtain \eqref{eq:ns52-support-cost}.
For the witness, $F_{f_a}(c+v)=\sqrt a F_h(av)$, so
$F_{f_a}(c)=0$, $\Psi_{f_a,f_a}'(c)=0$, and
$\Psi_{f_a,f_a}''(c)=2a^3M_1^2$. Apply the double-zero formula.
This is a lower bound on a conditional individual cluster's derivative,
not a lower bound on the complete arithmetic derivative: other clusters
remain present.
\end{proof}

\paragraph{Passage to the whole time-dependent sum.}
Absolute convergence of \eqref{eq:ns52-zero-form} at each time
does not imply convergence of its differentiated cluster series.
One sufficient hypothesis is an exhausting family of finite zero
multisets tracked with multiplicity on $[0,t_0]$, with piecewise
contours as necessary, whose endpoint sums converge to the complete
forms, and whose integrated derivatives satisfy
\[
 \int_0^{t_0}S_n'(t)[f]\,dt\le C\|f\|_2^2
\]
for every compact smooth $f$, with one $C$ independent of the
exhaustion and its support. An integrable bound
$S_n'(t)[f]\le B(t)\|f\|_2^2$, $\int_0^{t_0}B(t)\,dt<\infty$,
would suffice. Full symmetry makes these diagonal quantities real.
These are hypotheses, not consequences of pair cancellation. They
would imply $q_a[f]\ge q_{a,t_0}[f]-C\|f\|_2^2$ by the fundamental
theorem of calculus followed by the endpoint limit. The next result
shows that this particular comparison has an additional obstruction.

\begin{proposition}[An unmatched real zero obstructs uniform
same-test comparison]\label{prop:ns52-unmatched}
Assume the archived cutoff-radical lemma: there exists
$\phi\in\mathcal S(\mathbb R)$ such that, for every fixed finite
linear combination $h$ of its real translates, smooth cutoffs
$\chi_a h\in C_c^\infty(-a,a)$ satisfy
\[
 \chi_a h\longrightarrow h\text{ in }L^2\text{ and }L^1,
 \qquad q_a[\chi_a h]\longrightarrow0.
\]
Assume $q_{a,t_0}$ is the same-test zero form
\eqref{eq:ns52-zero-form}, all its zeros are real, and it has a
zero $c$ of multiplicity $m_c\ge1$ with $F_\phi(c)\ne0$.
Then no finite $C\ge0$ independent of $a$ can satisfy
\begin{equation}\label{eq:ns52-rejected-comparison}
 q_a[f]\ge q_{a,t_0}[f]-C\|f\|_2^2
 \quad(f\in C_c^\infty(-a,a))
\end{equation}
on a cofinal family of windows.
\end{proposition}
\begin{proof}
Fix $T>0$. The Schwartz assumption implies
\[
 B_T:=\sum_{k\in\mathbb Z}
       |\langle\phi,\phi(\cdot-kT)\rangle|<\infty.
\]
For each fixed positive integer $M$ put
\[
 h_M=M^{-1/2}\sum_{j=1}^M e^{icjT}\phi(\cdot-jT).
\]
Counting equal differences of shifts gives
$\|h_M\|_2^2\le B_T$, whereas
$F_{h_M}(c)=\sqrt M F_\phi(c)$.
Since all deformed zeros are real,
$q_{a,t_0}[f]\ge m_c|F_f(c)|^2$ for every smooth test.
Apply \eqref{eq:ns52-rejected-comparison} to $\chi_a h_M$ and
let $a\to\infty$ along the prescribed cofinal family, keeping $M$
fixed. The $L^1$ convergence controls the single real Fourier
evaluation. Thus
\[
 m_c M|F_\phi(c)|^2\le C\|h_M\|_2^2\le CB_T.
\]
Now let $M\to\infty$, a contradiction. No cutoff estimate uniform
in $M$, global closedness, or lower semicontinuity of a deformed
form is used.
\end{proof}

The cutoff input is Lemma~\ref{lem:v136-total-radicals}: its ordinary
norm convergence and vanishing complete operator residual imply the
displayed form limit. Schwartz decay and dominated convergence supply
the additional $L^1$ limit. The source used there has Fourier transform
a nonzero constant times $\Xi$, by \eqref{eq:v150-mellin} and the
source normalization, so $F_\phi(c)\ne0$ means precisely that this
real deformed zero is not an original Xi zero. An unmatched zero
at a specified $t_0$ has not been certified or proved in this
checkpoint. The proposition states its hypothesis explicitly; it
does not assert an unconditional arithmetic closure. It also does
not assert any negative direction of the original Weil form.

\paragraph{Disposition.}
The finite algebraic test succeeds: multiplicities and complete
pairing remove the inverse-gap singularity. What had to change is
the normalization, the analytic off-real sesquilinear product, and
retention of the exterior-cluster logarithmic derivative. What is
still missing is a complete signed transport estimate with support
independent error. Moreover the proposed same-test positive-time
comparison is impossible whenever there is an unmatched real
deformed zero. A broader heat approach would therefore need a
different comparison or a justified transport of tests that respects
the original radical; no such transport is supplied. These are
finite exact identities, proved implications with explicit hypotheses,
and a named open transport input. No new window, tail metric,
numerical certificate, G2 assertion, or RH assertion is made.
